\documentclass[12pt,a4paper]{article}

\usepackage[utf8]{inputenc}
\usepackage[T1]{fontenc}
\usepackage{lmodern}
\usepackage{amsmath,amssymb,amsthm}
\usepackage{enumitem}
\usepackage{xcolor}
\usepackage{geometry}
\geometry{a4paper, margin=2.5cm}
\usepackage{hyperref}

\hypersetup{
    colorlinks=true,
    linkcolor=blue!70!black,
    citecolor=green!50!black,
    urlcolor=blue!70!black
}

\theoremstyle{definition}
\newtheorem{definition}{Definition}[section]
\theoremstyle{plain}
\newtheorem{theorem}[definition]{Theorem}
\newtheorem{lemma}[definition]{Lemma}
\newtheorem{proposition}[definition]{Proposition}
\newtheorem{corollary}[definition]{Corollary}
\newtheorem{conjecture}[definition]{Conjecture}
\theoremstyle{remark}
\newtheorem{remark}[definition]{Remark}
\newtheorem{example}[definition]{Example}

\newcommand{\Var}{\mathrm{Var}}
\newcommand{\dist}{\mathrm{dist}}
\newcommand{\AC}{\mathrm{AC}}
\newcommand{\BV}{\mathrm{BV}}
\newcommand{\id}{\mathrm{id}}

\title{\textbf{Log-Distance Integrability for Dissipative Flows:\\
BV Selections on Fractal Attractors}}
\author{Lukas Geiger\\
\small Independent Researcher, Bernau im Schwarzwald, Germany\\
\small ORCID: \href{https://orcid.org/0009-0005-7296-1534}{0009-0005-7296-1534}%
\thanks{AI tools (Claude by Anthropic) were used for mathematical formalization, literature verification, and editorial refinement. All mathematical content was developed, verified, and approved by the author.}}
\date{March 2026}

\begin{document}
\maketitle

\begin{center}
\small
\textbf{MSC 2020:} 37L30, 47H09, 46B20, 35Q30, 26A45\\
\textbf{Keywords:} nearest-point projection, log-Lipschitz regularity,
bounded variation, attractor geometry, dissipative PDEs, Navier--Stokes
\end{center}

\begin{abstract}
We introduce a regularity framework for nearest-point projections
onto compact attractors of dissipative semiflows in Hilbert spaces.
The classical Lipschitz projection theorem (Federer, 1959) requires
positive reach---a condition incompatible with fractal attractor
geometry.  We replace it with a \emph{trajectory log-Lipschitz}
condition (TLL), which controls the projection modulus at the
geometrically natural scale of the attractor distance rather than
the partition mesh.  Combined with a \emph{log-distance
integrability} condition (LDI), this yields bounded-variation
selections along absolutely continuous trajectories---precisely the
regularity needed for Gronwall-type arguments in conditional
regularity criteria.  The framework applies to any dissipative
semiflow with a compact attractor of finite fractal dimension,
including the 3D Navier--Stokes equations, reaction-diffusion
systems, and damped wave equations.  Two verifiable sufficient
conditions for LDI are identified: bounded variation of the
log-distance function, and exponential sublevel time-measure decay.
\end{abstract}


% ==================================================================
\section{Introduction}
% ==================================================================

Let $H$ be a separable Hilbert space and $\mathcal{A} \subset H$ a
compact set---the global attractor of a dissipative semiflow
$\{S(t)\}_{t \ge 0}$.  A fundamental question in the qualitative
theory of dissipative PDEs is the regularity of the nearest-point
projection
\[
  P(x) \;\in\; \arg\min_{a \in \mathcal{A}} \|x - a\|_H
\]
as a function of time along a trajectory $u(t) = S(t)\,u_0$.  If the
map $t \mapsto P(u(t))$ is absolutely continuous (AC), Gronwall-type
energy estimates close automatically; if it is merely measurable,
critical error terms remain uncontrolled.

The classical theory (Federer~\cite{Federer1959}) guarantees Lipschitz
continuity of the nearest-point projection on the tubular neighbourhood
of sets with \emph{positive reach}.  However, attractors of dissipative
PDEs generically have non-integer fractal dimension and reach zero
(cf.~\cite{Temam1997}, Chapter~V).  This rules out the direct
application of Federer's theorem.

The purpose of this note is to introduce a weaker regularity class---
\emph{trajectory log-Lipschitz} projections---that is compatible with
fractal geometry and still suffices for bounded-variation (BV) control
of the projected trajectory.  Motivated by log-Lipschitz regularity
results for dissipative flows
(cf.~\cite{BaeCannone2016}), we observe that the naive
global log-Lipschitz condition
\begin{equation}\label{eq:global-LL-intro}
  \|P(x) - P(y)\|
  \;\le\;
  C\,\|x-y\|\,\bigl(1 + \log\tfrac{R}{\|x-y\|}\bigr)
\end{equation}
does \emph{not} imply BV of $P \circ u$ for AC curves~$u$: the
partition sums $\sum_i \Delta_i(1+\log(R/\Delta_i))$ diverge
logarithmically as the mesh refines (Section~\ref{sec:counterexample}).
The correct formulation replaces the increment-scale weight
$\log(R/\|x-y\|)$ with the attractor-distance weight
$\log(R/d(t))$, where $d(t) = \dist(u(t),\mathcal{A})$.  This
yields a trajectory-adapted condition whose BV sums converge
whenever a log-distance integrability condition (LDI) holds.


\paragraph{Main contributions.}
This paper makes three contributions: (i)~We introduce the \emph{trajectory log-Lipschitz} (TLL) condition and the \emph{log-distance integrability} (LDI) condition as the natural regularity scale for BV projections onto fractal attractors, strictly between Lipschitz (too strong) and H\"older (too weak). (ii)~We prove a BV chain rule (Theorem~\ref{thm:BV-chain}) showing that TLL\,+\,LDI suffices for bounded-variation selection regularity. (iii)~We identify sufficient geometric conditions (bounded variation of the log-distance, strong tangential cone condition) that are compatible with the fractal geometry of dissipative PDE attractors.

\begin{remark}[Input-output contract]
\label{rem:api-contract}
The central result of this paper can be stated as an input-output contract:
\begin{description}
\item[Input:] An absolutely continuous trajectory $u \in AC([0,T]; H)$, a compact attractor $\mathcal{A} \subset H$, and the nearest-point projection $P : \mathcal{A}^R \to \mathcal{A}$.
\item[Conditions:] (TLL) The projection $P$ is trajectory log-Lipschitz along $u$; (LDI) the log-distance weight $\log(R/d(t))$ is $\|u'\|$-integrable.
\item[Output:] The composed selection $v(t) := P(u(t))$ belongs to $BV([0,T]; H)$ with the explicit bound $\mathrm{Var}(v) \leq C_{\mathrm{TLL}} \int_0^T \|u'(t)\| \,\Omega(t)\,dt$.
\end{description}
This contract is designed to interface directly with the companion paper~\cite{FST-NS2026}: the output (BV selection) is precisely Assumption~$G'$ of that paper.

\smallskip\noindent
\textbf{Assumption~$G'$} (from~\cite{FST-NS2026}):
\emph{There exists a nearest-point selection
$v(t) = P(u(t)) \in \mathcal{A}$ such that
$v \in BV([0,T];\,H)$ for every $T < T^*$.}

\smallskip\noindent
Combined with Condition~D of~\cite{FST-NS2026} (an energy-type inequality for the attractor distance), Assumption~$G'$ implies global regularity for the 3D Navier--Stokes equations.  The attractor $\mathcal{A}$ may have fractal geometry---TLL and LDI are compatible with non-integer Hausdorff dimension.
\end{remark}


% ==================================================================
\section{Setup and Notation}
% ==================================================================

Throughout, $H$ denotes a separable real Hilbert space with norm
$\|\cdot\|$ and inner product $\langle\cdot,\cdot\rangle$.

\begin{definition}[Dissipative semiflow and attractor]
\label{def:semiflow}
A \emph{dissipative semiflow} on $H$ is a family
$\{S(t)\}_{t \ge 0}$ of continuous maps $S(t): H \to H$ satisfying
$S(0) = \id$, $S(t+s) = S(t) \circ S(s)$, and possessing a
\emph{compact global attractor}
$\mathcal{A} \subset H$: a compact, invariant
($S(t)\mathcal{A} = \mathcal{A}$ for all $t \ge 0$) set that
attracts every bounded subset of~$H$.
\end{definition}

\begin{definition}[Tubular neighbourhood and distance]
\label{def:tube}
For $R > 0$, the \emph{$R$-tubular neighbourhood} of $\mathcal{A}$ is
\[
  \mathcal{A}^R
  \;:=\;
  \{x \in H : \dist(x, \mathcal{A}) < R\},
\]
where $\dist(x, \mathcal{A}) := \inf_{a \in \mathcal{A}} \|x - a\|$.
For a curve $u: [0,T] \to \mathcal{A}^R$, we write
$d(t) := \dist(u(t), \mathcal{A})$.
\end{definition}

\begin{definition}[Nearest-point selection]
\label{def:selection}
A \emph{nearest-point selection} is a measurable map
$P: \mathcal{A}^R \to \mathcal{A}$ satisfying
$\|x - P(x)\| = \dist(x, \mathcal{A})$ for all
$x \in \mathcal{A}^R$, and $P|_{\mathcal{A}} = \id$.
Such a selection exists by the Kuratowski--Ryll-Nardzewski theorem
applied to the compact-valued multifunction
$x \mapsto \arg\min_{a \in \mathcal{A}} \|x-a\|$.
\end{definition}

\begin{proposition}[Minimal metric derivative selection]
\label{prop:minimal-selection}
Among the nearest-point selections
$P: \mathcal{A}^R \to \mathcal{A}$, there exists a
\emph{canonical selection} $u^*(t)$ along any absolutely continuous
trajectory $u \in \AC([0,T];\,H)$ defined by
\[
  u^*(t)
  \;:=\;
  \arg\min_{a \in \mathcal{A},\;\|u(t)-a\|=d(t)}
  |a|'(t),
  \qquad
  |a|'(t)
  \;:=\;
  \limsup_{h \to 0}
  \frac{\dist(a,\,P(u(t+h)))}{|h|},
\]
that minimises the upper metric derivative
among all nearest-point selections.  For a.e.\
$t \in [0,T]$ at which the nearest point is unique
(i.e., the set $\Pi(t) := \{a \in \mathcal{A} :
\|u(t)-a\| = d(t)\}$ is a singleton), the canonical selection
coincides with the unique nearest point and hence is
unambiguously defined.
\end{proposition}

\begin{proof}
The nearest-point set $\Pi(t) := \{a \in \mathcal{A} :
\|u(t)-a\| = d(t)\}$ is compact and non-empty for each~$t$.
The upper metric derivative functional $a \mapsto
\limsup_{h\to 0}\|P_{a}(u(t+h)) - a\|/|h|$ is lower
semicontinuous on the compact set~$\Pi(t)$, hence attains its
minimum.  This yields the existence of $u^*(t)$ for every~$t$.

\smallskip\noindent
\emph{Uniqueness.}\;
At times where $\Pi(t)$ is a singleton, the minimiser is
trivially unique.  We show that $\Pi(t)$ is a singleton for
a.e.~$t$.  The distance function $d(t) := \dist(u(t), \mathcal{A})$
is absolutely continuous (as the composition of the
$1$-Lipschitz map $\dist(\cdot, \mathcal{A})$ with the AC
curve~$u$).  An AC function is differentiable a.e.\ (Lebesgue).
At any $t$ with $d(t) > 0$ and $d'(t)$ existing, the nearest
point is unique: indeed, if two distinct points
$a_1, a_2 \in \mathcal{A}$ both realised $d(t)$, a standard
first-variation argument shows that the directional derivatives
of $\dist(\cdot, \mathcal{A})$ at $u(t)$ in direction
$(a_1 - a_2)/\|a_1 - a_2\|$ would differ in sign, contradicting
the differentiability of $d$ at~$t$
(cf.~\cite{PoliquinRockafellarThibault2000}; see
also~\cite{DeBlasiMyjak1993} and~\cite{Zamfirescu1990} for
the related $\sigma$-porosity of the ambiguity set in
uniformly convex Banach spaces).
On $\{t : d(t) = 0\}$ the nearest point is $u(t)$ itself.
Hence the nearest point is unique for a.e.~$t \in [0,T]$.
At the remaining times (a null set),
the selection among multiple nearest points is made by the
metric derivative minimisation and is measurable by
the Kuratowski--Ryll-Nardzewski theorem.
\end{proof}

\begin{remark}[Role of the canonical selection]
The canonical selection $u^*$ is the natural candidate for TLL
(Definition~\ref{def:TLL}): among all measurable nearest-point
selections, it has the smallest upper metric derivative and hence the
best chance of satisfying the log-Lipschitz bound.
Different selections may have different regularity (see the
Limitations subsection in Section~\ref{sec:discussion});
the minimal metric derivative selection is the ``best'' one.
\end{remark}

\begin{remark}[Uniqueness for non-convex sets]\label{rem:uniqueness-caveat}
For non-convex $\mathcal{A}$, the nearest-point set $\Pi(t)$
may contain more than one element, and the metric derivative
minimiser among them need not be unique.  The uniqueness
claimed in the proposition holds only at times where $\Pi(t)$
is a singleton---which is a.e.\ by the Zamfirescu--De~Blasi--Myjak
porosity result.  On the (null) set of ambiguous times, any measurable
selection among the minimisers of the metric derivative
suffices for the BV chain rule (Theorem~\ref{thm:BV-chain}),
since a change on a null set does not affect the integral
bound~\eqref{eq:BV-bound}.
\end{remark}


% ==================================================================
\section{Why Global Log-Lipschitz Fails}
\label{sec:counterexample}
% ==================================================================

We first show that the global log-Lipschitz condition
\eqref{eq:global-LL-intro}, even combined with log-distance
integrability, does not directly yield BV of $P \circ u$.

\begin{proposition}[Failure of the global LL bound to control variation]
\label{prop:divergence}
Let $P: \mathcal{A}^R \to \mathcal{A}$ satisfy
\eqref{eq:global-LL-intro} with constant $C > 0$.  Let
$u: [0,1] \to \mathcal{A}^R$ be Lipschitz with
$\|u'(t)\| = 1$ a.e.  For the uniform partition
$t_i = i/n$ ($i = 0, \ldots, n$), the upper bound obtained
from \eqref{eq:global-LL-intro} satisfies
\[
  \sum_{i=0}^{n-1} \|P(u(t_{i+1})) - P(u(t_i))\|
  \;\le\;
  C\,(1 + \log(nR)),
\]
which diverges as $n \to \infty$.  In particular, the global LL
condition does not provide a finite \emph{a priori} bound on
$\Var(P \circ u)$.
\end{proposition}

\begin{proof}
Each increment satisfies $\Delta_i := \|u(t_{i+1}) - u(t_i)\| = 1/n$.
By \eqref{eq:global-LL-intro}:
\[
  \sum_{i=0}^{n-1} \|P(u(t_{i+1})) - P(u(t_i))\|
  \;\le\;
  C \sum_{i=0}^{n-1} \frac{1}{n}\,
  \bigl(1 + \log(nR)\bigr)
  \;=\;
  C\,(1 + \log(nR)).
\]
As $n \to \infty$, this grows as $C\log n \to \infty$.
Since the bound $C(1 + \log(nR))$ holds for \emph{every}
partition with mesh $1/n$, the supremum over all such
partitions diverges: $\sup_n C(1 + \log(nR)) = \infty$.
Hence the global LL condition does not provide a finite
\emph{a priori} bound on $\Var(P \circ u)$ via the standard
estimate.  (Whether the variation itself is infinite depends
on the specific geometry of~$\mathcal{A}$; the point is that
the global LL condition is insufficient as a \emph{tool} for
establishing BV---no matter how the partition is chosen, the
LL-derived upper bound diverges with refinement.)
\end{proof}

\begin{remark}
The divergence is \emph{logarithmic}, not polynomial---much slower
than for H\"older maps (where $\sum \Delta_i^\alpha$ converges for
$\alpha < 1$).  This makes global LL ``almost'' sufficient for BV,
but the logarithmic accumulation over infinitely many partition
points defeats it.  The trajectory log-Lipschitz condition avoids
this by anchoring the log weight at the attractor distance $d(t)$,
which does not depend on the partition.
\end{remark}


% ==================================================================
\section{The Trajectory Log-Lipschitz Condition}
\label{sec:TLL}
% ==================================================================

\begin{definition}[Trajectory log-Lipschitz projection (TLL)]
\label{def:TLL}
Let $u \in \AC([0,T];\, H)$ with $u(t) \in \mathcal{A}^R$ for
all~$t$, and let $P: \mathcal{A}^R \to \mathcal{A}$ be a
nearest-point selection.  We say that $P$ satisfies the
\emph{trajectory log-Lipschitz condition} (TLL) along~$u$ if there
exist constants $C_{\mathrm{TLL}} > 0$ and $h_0 > 0$ such that for
all $t \in [0,T]$ with $d(t) > 0$ and all
$h \in \mathbb{R}$ with $0 < |h| < h_0$ and
$t+h \in [0,T]$
(i.e., both forward and backward increments):
\begin{equation}\label{eq:TLL}
  \|P(u(t+h)) - P(u(t))\|
  \;\le\;
  C_{\mathrm{TLL}}\,\|u(t+h) - u(t)\|\,
  \Omega(t),
  \tag{TLL}
\end{equation}
where $\Omega(t) := 1 + \log(R/d(t))$ is the
\emph{log-distance weight}.

On the set $\{t : d(t) = 0\}$ (where $u(t) \in \mathcal{A}$),
$P(u(t)) = u(t)$ by idempotence of~$P$.  For the transition
$d(t) = 0$, $d(t+h) > 0$, we have
$\|P(u(t+h)) - u(t)\| = \|P(u(t+h)) - P(u(t))\|$.
Since $u(t) \in \mathcal{A}$ is a candidate for the nearest point
of $u(t+h)$, the minimality of $P$ gives
$\|u(t+h) - P(u(t+h))\| \le \|u(t+h) - u(t)\|$.
Together with the triangle inequality:
\[
  \|P(u(t+h)) - u(t)\|
  \;\le\;
  \|P(u(t+h)) - u(t+h)\| + \|u(t+h) - u(t)\|
  \;\le\;
  2\,\|u(t+h) - u(t)\|.
\]
This gives $|v'|(t) \le 2\,|u'|(t)$ on $\{d = 0\}$.
\end{definition}

\begin{remark}[Convention $C_{\mathrm{TLL}} \ge 2$]
\label{rem:CTLL-convention}
The factor~$2$ on $\{d = 0\}$ is absorbed by the TLL bound
provided $C_{\mathrm{TLL}} \ge 2$ and $\Omega(t) \ge 1$.
Throughout, we adopt the convention
$C_{\mathrm{TLL}} := \max(C_{\mathrm{TLL}},\,2)$;
this changes the BV bound \eqref{eq:BV-bound} by at most a
factor of~$2$ and does not affect the integrability
of~\eqref{eq:LDI}.
\end{remark}

\begin{remark}[Geometric interpretation of TLL]
\label{rem:TLL-geometry}
The condition TLL asserts that the nearest-point projection,
evaluated along the trajectory~$u$, has a local Lipschitz constant
that blows up \emph{at most logarithmically} as the trajectory
approaches the attractor:
\[
  \mathrm{Lip}_{\mathrm{loc}}(P \circ u;\, t)
  \;\lesssim\;
  1 + \log\frac{R}{d(t)}.
\]
At distance $d(t) \sim R$ from the attractor, the projection is
essentially Lipschitz (the log term is $O(1)$).  As
$d(t) \to 0$, the log weight grows, reflecting the increasing
geometric complexity of the attractor boundary at fine scales.
The logarithmic rate is the \emph{critical threshold}: faster
growth (e.g., polynomial) would make LDI fail generically, while
slower growth (bounded) would imply Lipschitz projection---too
strong for fractal geometry.
\end{remark}

\begin{remark}[TLL as Aubin property of the projection correspondence]
\label{rem:TLL-Aubin}
The TLL condition admits a natural reformulation in the language of
variational analysis.  Define the \emph{projection correspondence}
\[
  \Gamma: t \;\mapsto\;
  \bigl\{a \in \mathcal{A} : \|a - u(t)\| = d(t)\bigr\}
\]
mapping each time to the set of nearest attractor elements.
The TLL condition (Definition~\ref{def:TLL}) is \emph{analogous}
to the \emph{Aubin property} (metric regularity) of the correspondence
$\Gamma$ at rate $1/|\log r|$.  Specifically, $\Gamma$ is
metrically regular at rate $\phi(r) = 1/|\log r|$ if
\[
  \dist\bigl(a,\, \Gamma(t')\bigr)
  \;\le\;
  L\,\phi\bigl(\|u(t') - u(t)\|\bigr)\,\|u(t') - u(t)\|
\]
for $a \in \Gamma(t)$ and $t'$ near~$t$.
Note that TLL and the Aubin property are \emph{not} equivalent in
general: TLL weights by $\Omega(t) = 1 + \log(R/d(t))$
(the \emph{attractor distance}), while the Aubin formulation
weights by $1/|\log\|u(t')-u(t)\||$ (the \emph{increment size}).
The two formulations coincide in the special case where
$d(t) \sim \|u(t')-u(t)\|$ (i.e., the trajectory increment and
the attractor distance are of the same order), which occurs for
trajectories approaching the attractor non-tangentially.
This analogy connects the projection regularity theory to the
Aubin criterion and metric regularity framework of
Rockafellar--Wets~\cite{RockafellarWets1998} and
Dontchev--Rockafellar~\cite{DontchevRockafellar2009}.
The key advantage is that Aubin-property criteria are
well-developed for set-valued maps and provide verifiable sufficient
conditions (e.g., via the Mordukhovich coderivative) that may be
applicable to the attractor projection.
\end{remark}

\begin{remark}[TLL vs global LL]
\label{rem:TLL-vs-LL}
The global log-Lipschitz condition \eqref{eq:global-LL-intro}
does \emph{not} imply TLL.  For $h \to 0$ at fixed~$t$:
\[
  \log\frac{R}{\|u(t+h)-u(t)\|}
  \;\approx\;
  \log\frac{R}{h\,|u'(t)|}
  \;=\;
  \log\frac{R}{|u'(t)|} + \log\frac{1}{h}
  \;\to\; \infty,
\]
while $\log(R/d(t))$ remains fixed.  The global condition
produces an \emph{increment-scale} log weight that diverges
with the partition mesh; TLL produces a \emph{distance-scale}
log weight that depends only on the current geometry.
This is the structural reason why TLL yields BV while
global LL does not.
\end{remark}


% ==================================================================
\section{Log-Distance Integrability}
\label{sec:LDI}
% ==================================================================

\begin{definition}[Log-distance integrability (LDI)]
\label{def:LDI}
In the setting of Definition~\ref{def:TLL}, we say that
\emph{log-distance integrability} holds if
\begin{equation}\label{eq:LDI}
  \int_{\{d>0\}} \|u'(t)\|\,\Omega(t)\,\mathrm{d}t
  \;=\;
  \int_{\{d>0\}} \|u'(t)\|\,
  \bigl(1 + \log\tfrac{R}{d(t)}\bigr)\,\mathrm{d}t
  \;<\; \infty.
  \tag{LDI}
\end{equation}
On the set $\{t : d(t) = 0\}$, $v(t) = u(t)$ and the BV
contribution is controlled by $\int_{\{d=0\}} \|u'(t)\|\,\mathrm{d}t
< \infty$ (since $u$ is AC), so only the region $\{d > 0\}$
is relevant for \eqref{eq:LDI}.  To state the BV bound
\eqref{eq:BV-bound} uniformly over all of $[0,T]$, we
extend $\Omega$ to $\{d=0\}$ by setting $\Omega(t) := 1$
there (consistent with the factor-$2$ bound of
Remark~\ref{rem:CTLL-convention}).
\end{definition}

\begin{remark}[Minimality of LDI]
\label{rem:LDI-minimal}
The condition LDI is positioned precisely between the known
regularity classes:
\begin{itemize}
  \item \textbf{Lipschitz projection} ($\Omega \equiv 1$):
    LDI reduces to $\int \|u'\| < \infty$, which is automatic
    for AC curves.  This is the Federer/positive-reach regime.
  \item \textbf{H\"older projection} ($\Omega(t) \sim d(t)^{-\beta}$,
    $\beta > 0$): LDI requires $\int \|u'\|\,d(t)^{-\beta} < \infty$,
    which generically fails near the attractor.  This is why
    Man\'e projectors (H\"older inverse) do not yield BV.
  \item \textbf{Log-Lipschitz projection} (TLL):
    LDI requires $\int \|u'\|\,\log(1/d) < \infty$, which is the
    borderline case---logarithmic singularities are integrable
    under mild dynamical conditions.
\end{itemize}
The log-Lipschitz/LDI pair thus occupies the \emph{critical}
regularity window for BV selections on fractal attractors.
\end{remark}


% ==================================================================
\section{Main Result: BV Chain Rule}
\label{sec:main}
% ==================================================================

\begin{theorem}[BV chain rule for log-Lipschitz projections]
\label{thm:BV-chain}
Let $\{S(t)\}_{t \ge 0}$ be a dissipative semiflow on a separable
Hilbert space~$H$ with compact global attractor
$\mathcal{A} \subset H$.  Let $u \in \AC([0,T];\, H)$ with
$u(t) \in \mathcal{A}^R$ for all~$t$, and let
$P: \mathcal{A}^R \to \mathcal{A}$ be a nearest-point selection
satisfying TLL \eqref{eq:TLL} along~$u$.  Set $v(t) := P(u(t))$.

If LDI \eqref{eq:LDI} holds, then
$v \in \BV([0,T];\, H)$ with the explicit bound
\begin{equation}\label{eq:BV-bound}
  \Var_H(v;\,[0,T])
  \;\le\;
  C_{\mathrm{TLL}}
  \int_0^T \|u'(t)\|\,\Omega(t)\,\mathrm{d}t,
\end{equation}
where $\Omega(t) := 1 + \log(R/d(t))$ on $\{d > 0\}$ and
$\Omega(t) := 1$ on $\{d = 0\}$
(Remark~\ref{rem:CTLL-convention}).
\end{theorem}

\begin{proof}
The proof has three steps.

\medskip
\noindent
\textbf{Step 1 (Metric derivative bound).}\;
For any continuous curve $v: [0,T] \to H$, the \emph{upper
metric derivative} at $t$ is
\[
  |v'|^+(t)
  \;:=\;
  \limsup_{h \to 0} \frac{\|v(t+h) - v(t)\|}{|h|}
  \;\in\; [0,+\infty].
\]
This quantity exists everywhere (as a value in $[0,+\infty]$)
and is Borel measurable.  For absolutely continuous curves the
$\limsup$ is a genuine limit a.e.\
(cf.~\cite[Theorem~1.1.2]{AmbrosioGigliSavare2008}).
Since $u$ is AC, $|u'|(t) := \lim_{h\to 0}\|u(t+h)-u(t)\|/|h|$
exists for a.e.~$t$.

\smallskip\noindent
\emph{Case $d(t)>0$:}\;
Dividing \eqref{eq:TLL} by $|h|$ and taking $\limsup$:
\begin{equation}\label{eq:metric-deriv-bound}
  |v'|^+(t)
  \;\le\;
  C_{\mathrm{TLL}}\,|u'|(t)\,\Omega(t)
  \qquad \text{for a.e.\ } t \text{ with } d(t) > 0.
\end{equation}
(The right-hand side is a genuine limit; the left-hand side
may be a $\limsup$, which is dominated by the same bound.)

\smallskip\noindent
\emph{Case $d(t)=0$ (interior of $\mathcal{Z}$):}\;
On $\mathcal{Z} := \{t \in [0,T] : d(t) = 0\}$,
$u(t) \in \mathcal{A}$ and $v(t) = P(u(t)) = u(t)$.
Since $u(t) \in \mathcal{A}$ is a candidate minimiser for
$\dist(u(t+h), \mathcal{A})$, the triangle inequality gives
$\|P(u(t+h)) - u(t)\| \le 2\,\|u(t+h) - u(t)\|$
(cf.\ the bound in Definition~\ref{def:TLL}),
so $|v'|^+(t) \le 2\,|u'|(t)$.

\smallskip\noindent
\emph{Boundary points of $\mathcal{Z}$:}\;
At $t_0 \in \partial\mathcal{Z}$ with $d(t_0)=0$, sequences
$t_n \to t_0$ with $d(t_n) > 0$ may exist.  For continuity of
$v$ at $t_0$: $\|v(t_n) - v(t_0)\| = \|P(u(t_n)) - u(t_0)\|
\le d(t_n) + \|u(t_n) - u(t_0)\| \to 0$, since $u$ is
continuous and $d = \dist(\cdot,\mathcal{A}) \circ u$ inherits
continuity (the distance function is $1$-Lipschitz).  The metric derivative bound
$|v'|^+(t_0) \le 2\,|u'|(t_0)$ remains valid by the same
triangle inequality argument as above.

\smallskip\noindent
Since $\Omega(t) \ge 1$ and $C_{\mathrm{TLL}} \ge 2$
(Remark~\ref{rem:CTLL-convention}), the
bound \eqref{eq:metric-deriv-bound} holds on all of $[0,T]$
(with $|v'|^+$ in place of $|v'|$ where the limit may not
exist; this suffices for the BV criterion in Step~3).

\medskip
\noindent
\textbf{Step 2 (Continuity of~$v$).}\;
The nearest-point projection on a compact (but non-convex) set
is in general only upper semicontinuous as a set-valued map
and need not be single-valued everywhere.  However, the
TLL condition \eqref{eq:TLL} directly implies continuity
of $v = P \circ u$ along the trajectory: for $d(t) > 0$,
$\|v(t+h) - v(t)\| \le C_{\mathrm{TLL}}\,\|u(t+h)-u(t)\|\,
\Omega(t) \to 0$ as $h \to 0$; for $d(t) = 0$,
$\|v(t+h)-v(t)\| \le 2\,\|u(t+h)-u(t)\| \to 0$
(including boundary points of $\mathcal{Z}$, cf.\ Step~1).
Hence $v$ is continuous on $[0,T]$.

\medskip
\noindent
\textbf{Step 3 (BV bound).}\;
We show $\|v(b) - v(a)\| \le \int_a^b g(t)\,\mathrm{d}t$
for all $0 \le a < b \le T$, where
$g(t) := C_{\mathrm{TLL}}\,|u'|(t)\,\Omega(t) \in L^1$
(by LDI).  Since $v$ is continuous (Step~2), the BV bound
$\Var(v) \le \int_0^T g$ then follows by summing over any
partition.

Set $f(t) := \|v(t) - v(a)\|$.  Then $f: [a,b] \to \mathbb{R}$
is continuous and real-valued,
$f(a) = 0$, and the upper Dini derivative satisfies
$D^+ f(t) \le |v'|^+(t)$ (since
$f(t+h) - f(t) \le \|v(t+h)-v(t)\|$ by the reverse triangle
inequality, and division by $h > 0$ followed by
$\limsup$ gives $D^+ f \le |v'|^+$).
By Step~1, $|v'|^+(t) \le g(t)$ a.e.
Moreover, $f(t) = \|v(t) - v(a)\| \ge 0$, so
$D_+ f(t) \ge -|v'|^+(t) > -\infty$ a.e.\ (since $g \in L^1$
implies $|v'|^+ < \infty$ a.e.).  A classical result on Dini
derivatives (see, e.g.,
\cite[Chapter~VII, \S\S\,9--10]{Saks1937}; cf.\ also
\cite[Chapter~3]{AmbrosioFuscoPallara2000}) gives:
for a continuous function $f:[a,b] \to \mathbb{R}$ with
$D^+ f(t) \le g(t)$ a.e., $D_+ f(t) > -\infty$ a.e., and
$g \in L^1$, the inequality
$f(b) - f(a) \le \int_a^b g(t)\,\mathrm{d}t$ holds.
Therefore
\[
  \|v(b) - v(a)\| \;=\; f(b)
  \;\le\; \int_a^b g(t)\,\mathrm{d}t
  \;\le\; \int_0^T g(t)\,\mathrm{d}t.
\]
Summing over any partition $0 = t_0 < \cdots < t_n = T$:
\[
  \sum_i \|v(t_{i+1}) - v(t_i)\|
  \;\le\;
  \sum_i \int_{t_i}^{t_{i+1}} g(t)\,\mathrm{d}t
  \;=\;
  \int_0^T g(t)\,\mathrm{d}t
  \;=\;
  C_{\mathrm{TLL}} \int_0^T |u'|(t)\,\Omega(t)\,\mathrm{d}t,
\]
which gives \eqref{eq:BV-bound}.

\medskip
\noindent
The bound \eqref{eq:BV-bound} follows from Steps~1--3.
\end{proof}

\begin{remark}[Direct partition-sum verification]
\label{rem:partition-sum}
For a fixed partition with mesh $< h_0$, the finiteness of the
partition sum can be seen directly from TLL without the
Dini-derivative machinery: on each subinterval, TLL gives
$\|v(t_{i+1}) - v(t_i)\|
\le C_{\mathrm{TLL}}\,\|u(t_{i+1}) - u(t_i)\|\,\Omega(t_i)$
(or the factor-$2$ bound on $\{d = 0\}$, absorbed by
$C_{\mathrm{TLL}} \ge 2$).  The resulting partition sum
$\sum_i C_{\mathrm{TLL}}\,\Omega(t_i)\int_{t_i}^{t_{i+1}}
\|u'\|\,\mathrm{d}t$ is finite for each fixed partition
(since $\Omega$ is bounded on any compact set disjoint
from $\{d = 0\}$, and the $\{d=0\}$ contribution is
controlled by $2\int\|u'\|$).
However, the passage to the supremum over \emph{all}
partitions---which is the definition of
$\Var(v)$---requires the Dini-derivative argument of
Step~3, which sharpens the pointwise $\Omega(t_i)$-weight
to the integrated $\Omega(t)$-weight.
\end{remark}

\begin{remark}[Banach space generality]
\label{rem:banach}
The proof of Theorem~\ref{thm:BV-chain} uses only the norm,
compactness of~$\mathcal{A}$, and measurable selection---no
inner product or Hilbert space structure.  The BV chain rule
therefore holds verbatim in any separable Banach space.  We
restrict to Hilbert spaces only for the PDE applications in
Section~\ref{sec:applications}.
\end{remark}


% ==================================================================
\section{Sufficient Conditions for LDI}
\label{sec:sufficient}
% ==================================================================

The condition LDI \eqref{eq:LDI} is a single integrability
requirement.  We provide two independent sufficient conditions,
each verifiable from dynamical information about the semiflow.

% ------------------------------------------------------------------
\subsection{Variant A: Log-distance has bounded variation}
% ------------------------------------------------------------------

\begin{proposition}[BV log-distance implies LDI]
\label{prop:variant-A}
Suppose $d(t) > 0$ for a.e.\ $t \in [0,T]$ and
$\log(R/d(t)) \in \BV([0,T])$.  Then LDI holds.
\end{proposition}

\begin{proof}
Write $\ell(t) := 1 + \log(R/d(t))$.  Since $u' \in L^1(0,T;\,H)$
(as $u$ is AC) and $\ell \in \BV([0,T]) \subset
L^\infty([0,T])$, the product
$\|u'(t)\|\,\ell(t)$ is locally integrable:
\[
  \int_0^T \|u'(t)\|\,\ell(t)\,\mathrm{d}t
  \;\le\;
  \|\ell\|_{L^\infty(0,T)}\,\int_0^T \|u'(t)\|\,\mathrm{d}t
  \;<\; \infty.  \qedhere
\]
\end{proof}

\begin{remark}[When is log-distance BV?]
\label{rem:log-BV}
The condition $\log(R/d) \in \BV$ requires that the trajectory
does not oscillate infinitely often between different distance
scales from the attractor.  Formally, it excludes infinite
``zickzack frequency'' in $d(t)$.  For monotonically attracting
trajectories ($d(t)$ eventually decreasing), this is automatic.
For trajectories with transient oscillations, BV of $\log d$
follows if the total number of scale crossings
$\{t : d(t) = \varepsilon\}$ is finite for each
$\varepsilon > 0$---a generic property of dissipative semiflows
with isolated equilibria.
\end{remark}

% ------------------------------------------------------------------
\subsection{Variant B: Sublevel time-measure control}
% ------------------------------------------------------------------

\begin{definition}[Sublevel-time control (STC)]
\label{def:STC}
The trajectory~$u$ satisfies \emph{sublevel-time control} with
exponent $\alpha > 0$ if there exists $C_* > 0$ such that
\begin{equation}\label{eq:STC}
  \bigl|\{t \in [0,T] : d(t) \le \varepsilon\}\bigr|
  \;\le\;
  C_*\,\varepsilon^\alpha
  \qquad \text{for all } \varepsilon \in (0, R].
  \tag{STC}
\end{equation}
\end{definition}

\begin{remark}[STC implies $|\{d=0\}| = 0$]
\label{rem:STC-d0}
Taking $\varepsilon \to 0$ in \eqref{eq:STC} shows
$|\{d = 0\}| = 0$.  Hence STC is only applicable to
trajectories that do not spend positive time on the attractor.
This is the generic case for non-stationary solutions
approaching~$\mathcal{A}$ asymptotically.  For trajectories
that reach the attractor in finite time (e.g., convergence
to an equilibrium with $u(t) \in \mathcal{A}$ for
$t \ge t_1$), Variant~A (BV log-distance) is the appropriate
sufficient condition.
\end{remark}

\begin{proposition}[STC implies LDI]
\label{prop:variant-B}
If STC \eqref{eq:STC} holds with some $\alpha > 0$ and
$u' \in L^2(0,T;\,H)$, then
$\log(R/d) \in L^p(0,T)$ for every $p < \infty$, and LDI holds.
\end{proposition}

\begin{proof}
\textbf{Step 1 (Tonelli representation).}\;
Since $u(t) \in \mathcal{A}^R$, we have $d(t) \le R$ for all~$t$.
On the set $\{t : d(t) = 0\}$ (where $u(t) \in \mathcal{A}$),
the integrand $\|u'(t)\|\,\log(R/d(t))$ is $+\infty$;
we set the contribution to zero by convention (on this set,
$v(t) = u(t)$ and the BV chain rule bound holds trivially).
On $\{t : d(t) > 0\}$, using the layer-cake identity
$\log(R/d(t)) = \int_{d(t)}^{R} s^{-1}\,\mathrm{d}s$ and Tonelli:
\begin{equation}\label{eq:tonelli}
  \int_{\{d>0\}} \|u'(t)\|\,\log\frac{R}{d(t)}\,\mathrm{d}t
  \;=\;
  \int_0^R \frac{1}{s}
  \Bigl(\int_{\{0 < d(t) \le s\}} \|u'(t)\|\,\mathrm{d}t\Bigr)
  \,\mathrm{d}s.
\end{equation}

\medskip
\noindent
\textbf{Step 2 (H\"older on inner integral).}\;
By H\"older's inequality:
\[
  \int_{\{0<d(t)\le s\}} \|u'(t)\|\,\mathrm{d}t
  \;\le\;
  \Bigl(\int_0^T \|u'(t)\|^2\,\mathrm{d}t\Bigr)^{1/2}
  \cdot
  \bigl|\{0 < d \le s\}\bigr|^{1/2}
  \;\le\;
  \|u'\|_{L^2}\,C_*^{1/2}\,s^{\alpha/2},
\]
where the last step uses $|\{0<d\le s\}| \le |\{d\le s\}|
\le C_*\,s^\alpha$ from STC.

\medskip
\noindent
\textbf{Step 3 (Integrability near $s=0$).}\;
Substituting into \eqref{eq:tonelli}:
\[
  \int_0^R \frac{1}{s}\,\|u'\|_{L^2}\,C_*^{1/2}\,s^{\alpha/2}
  \,\mathrm{d}s
  \;=\;
  \|u'\|_{L^2}\,C_*^{1/2}
  \int_0^R s^{\alpha/2 - 1}\,\mathrm{d}s
  \;=\;
  \|u'\|_{L^2}\,C_*^{1/2}\,
  \frac{2}{\alpha}\,R^{\alpha/2}
  \;<\; \infty.
\]
The integral converges because $\alpha > 0$, so
$s^{\alpha/2-1}$ is integrable near $s=0$.

\medskip
\noindent
\textbf{Step 4 ($L^p$ bounds).}\;
The same Tonelli argument with H\"older exponents $p$ and $q$
($1/p + 1/q = 1$) gives
$\|\log(R/d)\|_{L^p} \le C(p,\alpha,R,C_*)$ for every
$p < \infty$, so LDI follows from any finite $L^q$ bound
on~$\|u'\|$.
\end{proof}

\begin{remark}[Regularity of $u'$ for Navier--Stokes]
\label{rem:uprime-regularity}
Proposition~\ref{prop:variant-B} requires
$u' \in L^2(0,T;\,H)$.  For the 3D Navier--Stokes equations,
this holds for \emph{strong solutions}
($u' \in L^2(0,T;\,H)$ follows from $u \in L^2(0,T;\,V) \cap
L^\infty(0,T;\,H)$ and the equation).  For \emph{Leray--Hopf
weak solutions}, one only has
$u' \in L^{4/3}(0,T;\,V')$, where $V' = H^{-1}$.  In that
setting, the H\"older step (Step~2 of the proof) must be
adapted to use the $V'$-norm, with the STC exponent adjusted
accordingly.  The framework remains applicable; only the
quantitative constants change.
\end{remark}

\begin{remark}[Sources of STC in dissipative flows]
\label{rem:STC-sources}
Sublevel-time control \eqref{eq:STC} is the genuine new
analytical input.  Three mechanisms can generate it:

\begin{enumerate}[label=(\roman*)]
  \item \textbf{Exponential tracking:}\;
    If $d(t) \le C_0\,e^{-\lambda t}$ for $t \ge t_0$, then
    $\{d \le \varepsilon\} \subset
    \{t \ge (\log(C_0/\varepsilon))/\lambda\}$,
    whose measure is $O(\log(1/\varepsilon))$---much smaller than
    $C_*\,\varepsilon^\alpha$ for any $\alpha > 0$.
    So exponential tracking implies STC with \emph{any} exponent.
  \item \textbf{Lyapunov dissipation:}\;
    If a Lyapunov functional $E$ satisfies
    $-\dot{E} \ge \Psi(d)$ with $\Psi$ not too flat near $d=0$
    (e.g.\ $\Psi(d) \ge c\,d^\beta$ for some $\beta$),
    then sublevel times are controlled by the total energy budget
    $E(0) - \inf E$.
  \item \textbf{No fast oscillation:}\;
    Energetic constraints (finite dissipation integral)
    prevent arbitrarily rapid back-and-forth motion near the
    attractor, bounding the number of excursions into
    $\{d \le \varepsilon\}$.
\end{enumerate}
Each mechanism yields STC with different exponents~$\alpha$.
The STC condition is physically plausible for all standard
classes of dissipative PDEs.
\end{remark}

\begin{lemma}[Exponential tracking implies log-integrability]
\label{lem:stc-tracking}
If $d(t) \le C\,e^{-\lambda t}$ for some $\lambda > 0$ and
$C > 0$, then $\log(1/d) \in L^p([0,T])$ for all $p < \infty$
and all $T > 0$, and the LDI condition is satisfied with
\[
  \int_0^T \log\frac{1}{d(t)}\,\mathrm{d}t
  \;\le\;
  \frac{\lambda T^2}{2} + \log(1/C)\cdot T.
\]
In particular, STC holds with any exponent $\alpha > 0$, and
the LDI integral is bounded quadratically in~$T$.
\end{lemma}

\begin{proof}
From $d(t) \le C\,e^{-\lambda t}$ we obtain
$\log(1/d(t)) \ge \lambda t - \log C$, and also
\[
  \log\frac{1}{d(t)}
  \;\le\;
  \log\frac{1}{C\,e^{-\lambda t}}
  \;=\;
  \lambda t + \log\frac{1}{C}.
\]
Hence
\[
  \int_0^T \log\frac{1}{d(t)}\,\mathrm{d}t
  \;\le\;
  \int_0^T \bigl(\lambda t + \log(1/C)\bigr)\,\mathrm{d}t
  \;=\;
  \frac{\lambda T^2}{2} + \log(1/C)\cdot T.
\]
For the LDI condition \eqref{eq:LDI}, with $u' \in L^2$:
\[
  \int_0^T \|u'(t)\|\,\log\frac{R}{d(t)}\,\mathrm{d}t
  \;\le\;
  \|u'\|_{L^2}\,
  \Bigl(\int_0^T \bigl(\log(R/C) + \lambda t\bigr)^2\,\mathrm{d}t
  \Bigr)^{1/2}
  \;<\; \infty
\]
by H\"older's inequality, since the right factor is a polynomial
in~$T$.  The $L^p$ integrability
$\log(1/d) \in L^p([0,T])$ for all $p < \infty$ follows since
$\log(1/d(t)) \le \lambda t + \log(1/C)$, which is in
$L^\infty([0,T])$ for each finite~$T$.
\end{proof}

\begin{remark}[Exponential tracking as the strongest STC regime]
\label{rem:exp-tracking-strongest}
Lemma~\ref{lem:stc-tracking} shows that exponential tracking is
the ``easiest'' dynamical regime for LDI: the log-distance is
bounded by a linear function of time, making LDI automatic for
any trajectory with finite energy ($u' \in L^2$).
This is the regime expected for strong solutions of the
3D Navier--Stokes equations conditional on global regularity
(Example~\ref{ex:NS}).  The non-trivial case is when
exponential tracking fails---e.g., near potential blow-up---where
the weaker STC condition (polynomial sublevel decay) or
Variant~A (BV log-distance) become the relevant tools.
\end{remark}

\begin{remark}[STC as the realistic candidate for Navier--Stokes]
\label{rem:stc-realistic}
Of the three sufficient conditions for LDI (Variant~A: BV log-distance; Variant~B: sublevel time control; exponential tracking), \emph{sublevel time control} (STC) is the most realistic candidate for 3D Navier--Stokes:
\begin{enumerate}
\item \emph{Exponential tracking} ($d(t) \leq C_0 e^{-\lambda t}$) is the ``easy regime''---valid far from blow-up where the attractor is exponentially attracting. Near potential singularities, tracking rates may degenerate, making exponential tracking fragile.
\item \emph{STC} ($|\{d(t) \leq \varepsilon\}| \leq C_* \varepsilon^\alpha$) is robust: it only requires that the trajectory does not spend too long near the attractor at each scale. Two PDE mechanisms make STC plausible for Navier--Stokes:
\begin{itemize}
\item \emph{Dissipative time-scale separation:} Near the attractor, the low-mode dynamics are slow (controlled by the attractor) while the high-mode dynamics are fast (viscous damping). This separation limits the residence time at any given distance scale.
\item \emph{Finite determining modes:} The number of determining modes $N_d \sim G^{2/3}$ (where $G$ is the Grashof number) limits the effective dimension of the dynamics near $\mathcal{A}$, preventing pathological accumulation at any distance scale.
\end{itemize}
\item \emph{BV log-distance} (Variant~A) requires monotone approach to the attractor, which is generically violated by oscillatory solutions.
\end{enumerate}
\end{remark}


% ==================================================================
\section{Applications to Dissipative PDEs}
\label{sec:applications}
% ==================================================================

We illustrate how the abstract framework applies to concrete
PDE settings.  In each case, the semiflow possesses a compact
global attractor, and the question is whether TLL and LDI hold.

% ------------------------------------------------------------------
\subsection{3D Navier--Stokes on $\mathbb{T}^3$}
% ------------------------------------------------------------------

\begin{example}[Navier--Stokes equations]
\label{ex:NS}
Let $H = L^2(\mathbb{T}^3;\,\mathbb{R}^3)$ (divergence-free),
$\mathcal{A}$ the global attractor (in the sense of Temam~\cite{Temam1997}) with smooth
forcing~$f$.  By Temam~\cite{Temam1997}:
\begin{itemize}
  \item $\mathcal{A}$ is compact in $H^1$ (hence in $H$).
  \item Every $a \in \mathcal{A}$ is $C^\infty$.
  \item Exponential tracking holds \emph{conditional on global
    regularity}: for strong solutions existing on $[0,\infty)$,
    $d(t) \le C_0\,e^{-\lambda t}$ for $t \ge t_0$.
    (For Leray--Hopf weak solutions, only algebraic attraction
    in $L^2$ is known unconditionally.)
\end{itemize}

The attractor has finite fractal dimension
$d_f(\mathcal{A}) \le c_1\,G^{2/3}(1+\log G)^{1/3}$
(where $G$ is the Grashof number;
cf.~\cite{ConstantinFoiasTemam1985}), but generically has
reach~$0$---ruling out Federer's theorem.

\medskip\noindent
\textbf{Status of TLL:}\; Open.  The squeezing property of the
Navier--Stokes semiflow (cf.~\cite{FoiasManleyRosaTemam2001},
Chapter~14) provides exponential contraction of high-frequency
differences on the attractor.  Combined with Man\'e projector
estimates, this suggests that the nearest-point projection has
at most log-Lipschitz modulus along trajectories---but a rigorous
proof requires quantitative control of the projection stability
in the high-mode complement $Q_N H$.

\medskip\noindent
\textbf{Status of LDI:}\; Follows from exponential tracking
on $[t_0, T]$ for any fixed $T < \infty$
(Remark~\ref{rem:STC-sources}).  The critical question is the
blow-up regime $T \nearrow T^*$, where $\|u'\| \to \infty$.
The BV bound \eqref{eq:BV-bound} yields Assumption~G$'$ of
\cite{FST-NS2026}, which implies global regularity via the
attractor-distance argument.

\medskip\noindent
\textbf{Connection to the log-Dirichlet quotient:}\;
The companion paper~\cite{FST-NS2026} introduces the
\emph{log-Dirichlet quotient}
$Q_{\mathrm{LD}}(t) = \|\nabla(u-v)\|^2 / \|u-v\|^2$
and shows that its integrability transforms the Gronwall estimate
from exponential to algebraic growth.  The LDI condition of the
present paper provides precisely the mechanism that ensures
$Q_{\mathrm{LD}}$ integrability: the BV chain rule controls the
composition $\log d \circ u$, which bounds the log-Dirichlet quotient
from above.
\end{example}

% ------------------------------------------------------------------
\subsection{Reaction-diffusion systems}
% ------------------------------------------------------------------

\begin{example}[Reaction-diffusion on bounded domains]
\label{ex:RD}
Consider $\partial_t u = D\Delta u + f(u)$ on a bounded domain
$\Omega \subset \mathbb{R}^n$ with Dirichlet or Neumann
boundary conditions, where $D > 0$ is a diffusion matrix and
$f$ satisfies a dissipativity condition
$\langle f(u), u \rangle \le C - c|u|^2$.  The semiflow
possesses a compact global attractor $\mathcal{A} \subset L^2(\Omega)$
with finite fractal dimension~\cite{Temam1997}.

For scalar equations ($u \in \mathbb{R}$) in one space dimension
with a gradient nonlinearity $f(u) = -V'(u)$, the attractor
consists of equilibria and heteroclinic orbits, and the
nearest-point projection is Lipschitz (in fact, the attractor has
a Lipschitz-graph structure as an inertial manifold;
cf.~\cite{MalletParetSell1988}).  TLL and LDI hold trivially.

For systems in higher dimensions, the attractor may have non-integer
fractal dimension and reach zero.  The log-Lipschitz framework then
provides a path to BV selections that Federer's theorem does not.
The spectral gap condition for inertial manifolds is satisfied for
many reaction-diffusion systems~\cite{EdenFoiasNicolaenkoTemam1994},
making the attractor geometry more regular than in the
Navier--Stokes case.
\end{example}

% ------------------------------------------------------------------
\subsection{Damped wave equations}
% ------------------------------------------------------------------

\begin{example}[Damped semilinear wave equation]
\label{ex:wave}
Consider $\partial_{tt} u + \gamma\,\partial_t u = \Delta u + f(u)$
on $\Omega \subset \mathbb{R}^3$ with $\gamma > 0$ and subcritical
nonlinearity.  The semiflow on $H^1_0 \times L^2$ possesses a
compact global attractor with finite fractal dimension
(cf.~\cite{BabinVishik1992}).  The exponential tracking property
holds with rate $\lambda = \gamma/2$ (spectral gap of the damped
wave operator; cf.~\cite{BabinVishik1992}, Chapter~V).

The attractor geometry is intermediate between reaction-diffusion
(where inertial manifolds often exist) and Navier--Stokes (where
they do not).  The log-Lipschitz framework provides a
\emph{uniform} approach applicable across all three settings,
with TLL as the single geometric input.
\end{example}


% ==================================================================
\section{Discussion}
\label{sec:discussion}
% ==================================================================

\subsection{The new mathematical interface}

The central contribution of this note is the identification of
the \emph{critical regularity interface} for BV selections on
fractal attractors.  The interface is characterised by a single
question:

\begin{center}
\fbox{\parbox{0.85\textwidth}{%
\textbf{Core Question.}\;
For a dissipative semiflow with compact attractor
$\mathcal{A} \subset H$, does the nearest-point projection satisfy
trajectory log-Lipschitz control \eqref{eq:TLL}?
}}
\end{center}

\noindent
This is a question about the \emph{regularity of the attractor as
seen from the trajectory}, not about the global geometry
of~$\mathcal{A}$.  The distinction is crucial: global geometric
conditions (positive reach, inertial manifold) impose structure
on the \emph{entire} attractor, while TLL requires regularity
only in the direction of the dynamical approach.

\subsection{Relation to existing regularity theories}

\begin{center}
\begin{tabular}{lllll}
\hline
\textbf{Condition} & \textbf{Projection} & \textbf{Composition}
  & \textbf{Fractal?} & \textbf{Source} \\
\hline
Positive reach $> 0$
  & Lipschitz-1
  & AC $\Rightarrow$ AC
  & No
  & Federer~\cite{Federer1959} \\
Prox-regular
  & locally Lipschitz
  & AC $\Rightarrow$ AC (local)
  & No
  & Poliquin et al.~\cite{PoliquinRockafellarThibault2000} \\
Hausdorff-BV set map
  & BV selection exists
  & ---
  & Yes
  & Chistyakov~\cite{Chistyakov2004} \\
Man\'e projector
  & H\"older${}^{-1}$
  & AC $\Rightarrow$ H\"older
  & Yes
  & Man\'e~\cite{Mane1981} \\
\textbf{TLL + LDI}
  & \textbf{log-Lipschitz}
  & \textbf{AC $\Rightarrow$ BV}
  & \textbf{Yes}
  & \textbf{This note} \\
\hline
\end{tabular}
\end{center}

\noindent
The TLL+LDI framework is the only entry that simultaneously
(i)~accommodates fractal geometry (non-integer dimension, reach~$0$),
(ii)~yields BV regularity (sufficient for Gronwall-type arguments),
and (iii)~is formulated as a trajectory-level condition
(independent of global attractor structure).

\subsection{Limitations}

\begin{itemize}
  \item \textbf{Selection dependence.}\;
    The TLL condition is formulated for a \emph{given} measurable
    nearest-point selection~$P$.  On non-convex compact sets, the
    nearest-point map is generically multi-valued, and different
    selections may have different regularity.  The framework requires
    the existence of \emph{at least one} selection satisfying TLL;
    it does not claim that every measurable selection does so.  In
    practice, the ``best'' selection (e.g., the one minimising the
    metric derivative) is the natural candidate.
  \item \textbf{TLL is open for 3D NS.}\;
    The squeezing property and Man\'e projector estimates suggest
    log-Lipschitz stability of the projection, but a rigorous proof
    requires quantitative bounds on the high-mode projection
    complement that are not currently available.
  \item \textbf{LDI is non-trivial only under blow-up.}\;
    For smooth solutions on compact time intervals, LDI holds
    automatically.  The condition becomes critical only in the
    limit $T \nearrow T^*$ (blow-up time), where $\|u'\| \to \infty$.
  \item \textbf{BV, not AC.}\;
    The framework yields BV regularity, which suffices for
    Stieltjes-integral-based Gronwall arguments but not for
    pointwise differential inequalities.  Upgrading BV to AC
    would require Lipschitz projection (positive reach)---exactly
    the condition we are trying to avoid.
  \item \textbf{BV bound depends on $\|u'\|$.}\;
    The bound \eqref{eq:BV-bound} is
    $\Var(v) \le C_{\mathrm{TLL}} \int_0^T \|u'\|\,\Omega\,dt$.
    For Navier--Stokes, $\|u'\|$ depends on the solution norm,
    and near a potential blow-up ($T \nearrow T^*$),
    $\|u'\| \to \infty$.  The Gronwall closure requires that the
    BV bound can be ``absorbed'' before $\Var(v)$ diverges.  The
    details of this absorption depend on the companion
    paper~\cite{FST-NS2026} and are beyond the scope of this note.
\end{itemize}


\subsection*{Central open problem}
\begin{quote}
\textbf{Selection Regularity Problem.}
Let $u(t)$ be a Leray--Hopf solution of the 3D Navier--Stokes equations and $\mathcal{A}$ the global attractor. Does the nearest-point projection $P_\mathcal{A}$ satisfy the TLL condition along $u(t)$? Equivalently: does the attractor geometry permit BV-regular selections?

An affirmative answer, combined with the BV chain rule of this paper, would yield Assumption~G of the companion paper \cite{FST-NS2026} and hence global regularity.
\end{quote}


% ==================================================================
\section{Conclusion}
% ==================================================================

We have introduced a regularity framework for nearest-point
projections onto compact attractors that occupies the critical
gap between H\"older (insufficient for BV) and Lipschitz
(incompatible with fractal geometry).  The trajectory log-Lipschitz
condition (TLL) controls the projection modulus at the attractor
distance scale, and the log-distance integrability condition (LDI)
ensures that the resulting BV sums converge.

The framework reduces conditional regularity criteria for
dissipative PDEs---including the 3D Navier--Stokes
equations---to a single geometric-dynamical question: does the
nearest-point projection have at most logarithmic blowup along
dissipative trajectories?  This question is independent of the
specific PDE and depends only on the interplay between attractor
geometry and the semiflow's contraction properties.


% ==================================================================
\section{Time-averaged selection and BV regularity (speculative)}
\label{sec:time-averaged}
% ==================================================================

\emph{This section is speculative: it outlines an alternative approach
whose key steps remain open conjectures.  The results of
Sections~\ref{sec:main}--\ref{sec:applications} do not depend on
this material.}

\medskip

The central obstruction to closing the regularity gap is the BV
control of the nearest-point selection
$v(t) \in \arg\min_{a \in \mathcal{A}} \|u(t) - a\|$.  We propose
replacing the pointwise selection by a \emph{time-averaged} selection
that is automatically BV-regular.

\begin{definition}[Moreau--Yosida regularised selection]
\label{def:my-selection}
For $\varepsilon > 0$, define the \emph{time-averaged nearest-point
selection}
\[
  v_\varepsilon(t)
  \;:=\;
  \arg\min_{a \in \mathcal{A}}
  \int_{t-\varepsilon}^{t+\varepsilon}
  \|u(s) - a\|^2 \, ds.
\]
\end{definition}

\begin{conjecture}[Properties of $v_\varepsilon$]
\label{conj:bv-regularity}
Under the standing assumptions (finite-dimensional attractor, energy
dissipation bound), the regularised selection~$v_\varepsilon$ is expected to satisfy:
\begin{enumerate}[label=(\roman*)]
\item \emph{Uniqueness:}  $v_\varepsilon(t)$ is unique for a.e.~$t$
  (strict convexity of the time-integrated cost on~$H$; for minimisation
  over the non-convex compact set~$\mathcal{A}$, uniqueness requires
  a separate argument via Hausdorff stability of the minimiser set,
  which remains open).
\item \emph{BV bound:}
  $\Var(v_\varepsilon;\,[0,T])
   \;\le\;
   C_\varepsilon
   \int_0^T \|\partial_t u(s)\|\,ds
   \;<\; \infty$,
  where $C_\varepsilon = T/(2\varepsilon) \to \infty$ as
  $\varepsilon \to 0$.  A \emph{uniform} BV bound (independent
  of~$\varepsilon$) would require additional geometric input;
  see Remark~\ref{rem:epsilon-limit}.
\item \emph{TLL inheritance:}  If~$u$ satisfies TLL with constant
  $C_{\mathrm{TLL}}$, then $v_\varepsilon$ is expected to satisfy TLL with
  constant $C_{\mathrm{TLL}} + O(\varepsilon)$.  A rigorous proof
  requires control of the mismatch between pointwise and averaged
  distance, which remains open.
\item \emph{Gronwall compatibility:}  The Gronwall estimates of the
  main theorem (Theorem~\ref{thm:BV-chain}) are expected to remain valid
  for~$v_\varepsilon$ with an additional error term
  $O(\varepsilon\,\|\partial_t u\|_{L^2})$.  A rigorous formulation
  requires specifying the relative entropy functional~$H$ and the
  constants $\alpha, \beta, \gamma$; this is deferred to future work.
\end{enumerate}
\end{conjecture}

\begin{proof}[Partial justification]
\emph{(i)}\; The map
$a \mapsto \int_{t-\varepsilon}^{t+\varepsilon} \|u(s)-a\|^2\,ds$
is strictly convex on~$H$.  However, the minimisation is over the
compact, generally \emph{non-convex} set~$\mathcal{A}$.  Uniqueness
of the minimiser over a non-convex set requires a separate argument
(e.g., via Hausdorff stability of the minimiser set or a generic
transversality argument).  This step remains open for general
attractors; it holds when~$\mathcal{A}$ is a Lipschitz manifold.

\emph{(ii)}\; The first-variation condition for the minimiser yields
\[
  \|v_\varepsilon(t+h) - v_\varepsilon(t)\|
  \;\le\;
  \frac{1}{2\varepsilon}
  \int_{t-\varepsilon}^{t+\varepsilon}
  \|u(s+h) - u(s)\|\,ds
  \;\le\;
  \frac{h}{2\varepsilon}
  \int_{t-\varepsilon}^{t+\varepsilon}
  \|\partial_t u(s)\|\,ds.
\]
Integrating over $t \in [0,T]$ and applying Fubini gives the BV bound
with $C_\varepsilon = T/(2\varepsilon)$.
(Note: $C_\varepsilon \to \infty$ as $\varepsilon \to 0$; the bound is
uniform in~$t$ for each fixed $\varepsilon > 0$, but the passage to the
limit $\varepsilon \to 0$ requires an independent uniform BV bound---see
Remark~\ref{rem:epsilon-limit}.)

\emph{(iii)}\; The TLL constant degrades by the time-averaging window:
\[
  |\log d_\varepsilon(t+h) - \log d_\varepsilon(t)|
  \;\le\;
  C_{\mathrm{TLL}}\,|\log|h||
  \;+\;
  C\,\varepsilon/|h|,
\]
where the second term arises from the mismatch between pointwise and
averaged distance.

\emph{(iv)}\; The Gronwall estimate picks up an additional forcing
term from the time-averaging:
\[
  \frac{d}{dt} H(u \mid v_\varepsilon)
  \;\le\;
  -\alpha\,H + \beta + \gamma\,\varepsilon,
\]
where $\gamma$ depends on $\|\partial_t u\|$ but not on~$H$.
The standard Gronwall lemma applies; the additional
$\gamma\varepsilon$ term vanishes in the limit
$\varepsilon \to 0$.
\end{proof}

\begin{remark}[Passage to the limit $\varepsilon \to 0$]
\label{rem:epsilon-limit}
The key question is whether the BV bound survives the limit
$\varepsilon \to 0$:
\begin{itemize}
\item If\; $\sup_\varepsilon \Var(v_\varepsilon;\,[0,T]) < \infty$,
  then a subsequence $v_{\varepsilon_k}$ converges in~$L^1$ to a BV
  selection~$v_0$, which is the \emph{canonical} nearest-point
  selection.
\item The uniform BV bound holds if and only if the attractor
  $\mathcal{A}$ has a ``time-averaged reach''
  $\mathrm{reach}_{\mathrm{avg}}(\mathcal{A}) > 0$, which is
  strictly weaker than positive reach: it allows isolated cusps of
  measure zero.
\end{itemize}
This is the precise formulation of the ``BV regularity gap'' and
connects to:
\begin{itemize}
\item \emph{Moreau--Yosida regularisation} in convex analysis (the
  time-averaged cost is the Moreau envelope),
\item \emph{Viscous selection} in gradient systems (the time average
  acts as artificial viscosity),
\item \emph{Thermodynamic selection} in the Dark Energy companion
  paper (the IR averaging of~$\Phi[\rho]$).
\end{itemize}
\end{remark}


% ==================================================================
\section{Post-Zookeeper Transfer: From Proof of Life to Damped 3D Navier--Stokes}
\label{sec:zookeeper-ldi-transfer}

The Zookeeper normal form makes this paper the natural test carrier for
the Navier--Stokes branch.  Here the defect, gap, and selection mechanism
are already explicit:

In the updated CORE reading, this statement means RFEP v1.8 imports a
normal-form checklist from the Zookeeper closure of Riemann $C^{(2)}$/MS2.
It is not yet a proof of three-dimensional Navier--Stokes regularity; the
damped-to-classical transfer remains the branch-specific obligation.
\begin{description}
  \item[Operator/form.] The distance-to-attractor map and its measurable
  nearest-point selections.
  \item[Defect.] The log-distance integral, projection variation, and BV
  defect.
  \item[Approximation.] Lorenz $\to$ Kuramoto--Sivashinsky $\to$
  reaction-diffusion $\to$ damped 3D Navier--Stokes $\to$ classical
  3D Navier--Stokes.
  \item[Gap.] TLL, squeezing, exponential tracking, and damping gaps.
  \item[Limit.] From finite-dimensional attractors to PDE attractors, and
  from damped/hyperviscous systems to the classical $\beta=1$ equation.
\end{description}

The immediate next target is therefore not the full Millennium problem.
It is a damped 3D Navier--Stokes transfer theorem: in a parameter regime
where the dynamical-systems literature supplies a global attractor and
improved regularity, prove that damping plus attractor regularity implies
TLL/LDI and hence BV selection.  The existing Lorenz and
Kuramoto--Sivashinsky tests then become the first two rungs of a larger
ladder rather than isolated numerical evidence.

\section{Open Directions}
\label{sec:open-directions}
% ==================================================================

The TLL+LDI framework opens several concrete research directions
grouped by theme.

% ------------------------------------------------------------------
\subsection{Pilot verification systems for TLL and LDI}
% ------------------------------------------------------------------

\begin{remark}[Pilot verification systems for TLL and LDI]
\label{rem:pilot-verification}
Three candidate systems for pilot verification of the TLL and LDI conditions:
\begin{enumerate}
\item \emph{Reaction-diffusion (2D):} Smooth nonlinearity, well-understood attractors. TLL and LDI are expected to hold due to the parabolic smoothing.
\item \emph{Hyperviscous NSE ($(-\Delta)^\beta$, $\beta \geq 5/4$):} Inertial manifolds exist (Mallet-Paret--Sell, 1988), guaranteeing TLL. Serves as a bridge system between the tractable reaction-diffusion case and full NS ($\beta = 1$).
\item \emph{Kuramoto--Sivashinsky (1D):} Deterministic chaos without singularities. The attractor is finite-dimensional and smooth, making TLL+LDI verifiable by direct computation.
\end{enumerate}
\end{remark}

\begin{remark}[Proof of life: TLL and LDI on the Lorenz attractor]
\label{rem:tll-ldi-lorenz}
We verify TLL and LDI numerically on the Lorenz attractor ($\sigma = 10$, $\rho = 28$, $\beta = 8/3$), a 3D chaotic system with fractal attractor dimension $\approx 2.06$. The attractor is approximated by $50{,}000$ trajectory points; a test trajectory is initialised with perturbation $(2, 2, 2)$ from an attractor element and integrated for $T = 100$ time units.

\emph{Results} (all five diagnostics passed):

\begin{center}
\small
\begin{tabular}{l|l|l}
\textbf{Test} & \textbf{Result} & \textbf{Detail} \\
\hline
TLL (Definition~\ref{def:TLL}) & PASS & $C_{\mathrm{TLL}}(95\%) = 0.19$, max $= 1.94$ \\
LDI (Definition~\ref{def:LDI}) & PASS & $\int \|u'\| \Omega\,\mathrm{d}t = 6.47 \times 10^4$ (finite) \\
BV Chain Rule (Theorem~\ref{thm:BV-chain}) & PASS & $\mathrm{Var}(v) / (C_{\mathrm{TLL}} \cdot \mathrm{LDI}) = 0.52$ \\
STC (Definition~\ref{def:STC}) & PASS & $\alpha = 1.11 > 0$ \\
Exponential tracking (Lemma~\ref{lem:stc-tracking}) & PASS & $\lambda = 0.064 > 0$
\end{tabular}
\end{center}

Key observations: (i) the TLL constant at the 95th percentile
is $C_{\mathrm{TLL}} \approx 0.19$, but the \emph{maximum} over
the test trajectory is $1.94$---an order of magnitude larger.
Since TLL must hold for \emph{all} $t$ (not 95\%), the relevant
constant for the BV bound is the maximum.
(ii) The LDI/L$^1$ ratio is $6.86$, showing that the log-distance
weight inflates the integral by less than one order of magnitude.
(iii) The BV chain rule bound is satisfied with margin (ratio $0.52 < 1$).
(iv) Sublevel-time control holds with exponent $\alpha \approx 1.1$.

\emph{Limitations of the numerical test:}
The attractor approximation uses $50{,}000$ trajectory points, which
systematically overestimates $d(t)$ relative to the true attractor
(making TLL constants appear smaller).  Moreover, this is a single
test at the classical Lorenz parameters ($\sigma=10$, $\rho=28$,
$\beta=8/3$) in $\mathbb{R}^3$---a finite-dimensional ODE, not an
infinite-dimensional PDE.  A systematic parameter scan over $\rho$
(including near-onset values $\rho \approx 24$) would strengthen
the evidence.  Despite these caveats, the test confirms that the
TLL+LDI conditions are not vacuous: they are satisfied in at least
one concrete chaotic dynamical system.
\end{remark}

\begin{remark}[Proof of life~II: TLL and LDI on the Kuramoto--Sivashinsky attractor]
\label{rem:tll-ldi-ks}
We extend the numerical validation to the Kuramoto--Sivashinsky (KS) equation
$u_t + u\,u_x + u_{xx} + u_{xxxx} = 0$ on $[0, L]$ with periodic boundary
conditions ($L = 32\pi$), a 1D dissipative PDE with spatiotemporal chaos and
finite-dimensional attractor.  The pseudo-spectral discretisation uses $N = 128$
Fourier modes; the attractor is approximated by $5{,}001$ points in the
$M = 32$-dimensional Fourier coefficient space.

\emph{Results} (primary run, $N = 128$):

\begin{center}
\small
\begin{tabular}{l|l|l}
\textbf{Test} & \textbf{Result} & \textbf{Detail} \\
\hline
TLL (Definition~\ref{def:TLL}) & PASS & $C_{\mathrm{TLL}}(95\%) = 9.0$, max $= 58.8$ \\
LDI (Definition~\ref{def:LDI}) & PASS & $\int \|u'\| \Omega\,\mathrm{d}t = 5576$ (finite), LDI/$L^1$ $= 1.86$ \\
BV Chain Rule (Theorem~\ref{thm:BV-chain}) & PASS & $\mathrm{Var}(v) / (C_{\mathrm{TLL}} \cdot \mathrm{LDI}) = 0.059$ \\
STC (Definition~\ref{def:STC}) & PASS & $\alpha = 0.91 > 0$ \\
Exponential tracking & N/A & $\lambda = -0.074$ (chaotic wandering, no monotone approach)
\end{tabular}
\end{center}

\emph{Grid refinement} ($N = 32, 64, 128, 256$; $N=32$ uses $L = 22$):

\begin{center}
\small
\begin{tabular}{r|r|r|r|r|r|c}
$N$ & $M$ & $C_{\mathrm{TLL}}(95\%)$ & LDI/$L^1$ & BV ratio & STC $\alpha$ & All \\
\hline
32  & 16 & 1.81 & 2.58 & 0.145 & 1.85 & PASS \\
64  & 32 & 8.75 & 1.41 & 0.083 & 1.20 & PASS \\
128 & 32 & 5.10 & 2.02 & 0.078 & 1.11 & PASS \\
256 & 32 & 2.66 & 2.77 & 0.099 & 0.79 & PASS
\end{tabular}
\end{center}

Key observations: (i)~All four diagnostics pass at every resolution---the
TLL/LDI framework is \emph{not} an artefact of finite discretisation.
(ii)~The BV ratio is uniformly small ($< 0.15$), indicating the chain rule
bound is far from tight.
(iii)~Unlike the Lorenz system ($\mathbb{R}^3$), KS is a \emph{PDE}
with infinite-dimensional phase space, truncated to $N$ Fourier modes.
The attractor dimension ($\sim 8$--$12$ for $L = 32\pi$) is well above
the Lorenz value ($\sim 2.06$), making this a substantially harder test.
(iv)~Exponential tracking fails (expected: the test trajectory wanders
\emph{on} the attractor, not toward it).
Script: \texttt{compute\_tll\_ldi\_ks.py}; results: \texttt{ks\_tll\_ldi\_results.json}.
\end{remark}

\begin{remark}[Ginzburg--Landau as TLL/LDI bridge]
\label{rem:ginzburg-landau}
The complex Ginzburg--Landau equation $\partial_t u = (1+i\alpha)\Delta u + u - (1+i\beta)|u|^2 u$ possesses a finite-dimensional inertial manifold for certain parameter ranges. Complete verification of TLL+LDI on this system, followed by a lifting argument to 2D NS (via parameter continuation in the nonlinearity), would provide a concrete path to the full NS regularity problem.
\end{remark}

\begin{theorem}[Squeezing implies TLL for finite-dimensional attractors]\label{thm:squeezing-tll}
Let $u_t + Au + B(u,u) = f$ be a dissipative PDE on a Hilbert space $H$
with compact global attractor $\mathcal{A}$ of finite fractal dimension $d_F$.
Suppose the squeezing property holds: there exist $N \in \mathbb{N}$,
$\theta \in (0,1)$ and $T > 0$ such that for all $u_0, v_0$ in the
absorbing ball,
\[
  \|Q_N(S(T)u_0 - S(T)v_0)\| \leq \theta\,\|u_0 - v_0\|
  \quad\text{or}\quad
  \|Q_N(S(T)u_0 - S(T)v_0)\| \leq \|P_N(S(T)u_0 - S(T)v_0)\|,
\]
where $P_N$ projects onto the first $N$ eigenmodes of $A$.
Then the trajectory log-Lipschitz condition (TLL, Definition~\ref{def:TLL})
holds with
\[
  C_{\mathrm{TLL}} \leq C(N, d_F, \theta)\,\bigl(1 + |\log\theta|\bigr).
\]
\end{theorem}

\begin{proof}
The argument proceeds in three steps.

\emph{Step~1 (Lipschitz graph).}
The squeezing property implies that the global attractor $\mathcal{A}$
lies in a Lipschitz graph over the finite-dimensional subspace $P_N H$.
This is the Mané projection theorem applied to the squeezing dichotomy:
since the high-mode component $Q_N$ contracts relative to the low-mode
component $P_N$, the restriction $P_N|_{\mathcal{A}}$ is injective with
Lipschitz inverse \cite{FoiasTemam1979,Mane1981}.
Concretely, there exists $L = L(N, d_F) < \infty$ such that
$\|Q_N(u - v)\| \leq L\,\|P_N(u - v)\|$ for all $u, v \in \mathcal{A}$.

\emph{Step~2 (Nearest-point projection is locally Lipschitz on the graph).}
We must show that the \emph{nearest-point projection} $\pi_{\mathcal{A}}$
(not the linear projection $P_N$) is locally Lipschitz in a tubular
neighbourhood of the Lipschitz graph $\mathcal{A} = \mathrm{graph}(\Phi)$
where $\Phi: P_N H \to Q_N H$ with $\mathrm{Lip}(\Phi) \leq L$.

A Lipschitz graph with constant $L$ is \emph{prox-regular} in the sense
of Poliquin--Rockafellar--Thibault~\cite{PoliquinRockafellarThibault2000}:
the reach satisfies $\mathrm{reach}(\mathrm{graph}(\Phi)) \geq r_0 > 0$
(depending on $L$ and the curvature of the embedding), and on the
$r_0$-tubular neighbourhood the nearest-point projection
$\pi_{\mathcal{A}}$ is single-valued and Lipschitz.  More precisely,
for any $x, y$ with $\dist(x, \mathcal{A}) < r_0$ and
$\dist(y, \mathcal{A}) < r_0$,
\[
  \|\pi_{\mathcal{A}}(x) - \pi_{\mathcal{A}}(y)\|
  \;\leq\;
  \frac{\sqrt{1 + L^2}}{1 - \max(\dist(x,\mathcal{A}), \dist(y,\mathcal{A}))/r_0}
  \,\|x - y\|.
\]
For points $u(t)$ on the trajectory with $d(t) = \dist(u(t), \mathcal{A})$,
the denominator stays bounded away from zero as long as $d(t) < r_0/2$
(which holds eventually by the absorbing ball property).  This gives
$\|\pi_{\mathcal{A}}(u(t)) - \pi_{\mathcal{A}}(u(t+h))\|
\leq 2\sqrt{1+L^2}\,\|u(t+h) - u(t)\|$
for $u(t), u(t+h)$ in the $(r_0/2)$-tube.

\emph{Step~3 (Log factor from exponential approach).}
For trajectories $u(t)$ in the absorbing ball but not yet on $\mathcal{A}$,
the squeezing property gives exponential approach:
$d(u(t), \mathcal{A}) \leq C\,\theta^{t/T}\,d(u(0), \mathcal{A})$.
Taking logarithms, $\log(R/d(u(t), \mathcal{A})) \geq (t/T)\,|\log\theta| - C'$.
The TLL modulus $\Omega(t) = 1 + \log(R/d(t))$ therefore grows at most linearly
in $t$, while the exponential contraction in the high modes provides the
compensating decay.  Combining with Steps~1--2 and the exponential tracking
estimates of Eden--Foias--Nicolaenko--Temam \cite{EdenFoiasNicolaenkoTemam1994}, one obtains
the bound $C_{\mathrm{TLL}} \leq C(N, d_F, \theta)\,(1 + |\log\theta|)$.
\end{proof}

\begin{remark}[Hyperviscous Navier--Stokes as complete TLL/LDI test case]
\label{rem:hyperviscous}
For the hyperviscous NS equation $\partial_t u + (-\Delta)^\beta u + (u\cdot\nabla)u = -\nabla p + f$
with $\beta \geq 5/4$ on $\mathbb{T}^3$, inertial manifolds exist
(Mallet-Paret--Sell~\cite{MalletParetSell1988}, Romanov~\cite{Romanov2000}). The squeezing property holds
with explicit constants, so Theorem~\ref{thm:squeezing-tll} applies directly:
TLL holds with $C_{\mathrm{TLL}} \leq C(\beta, \nu, f)$.
Combined with LDI (which follows from the exponential tracking on the
inertial manifold), this gives a \emph{complete closed cycle}:
\[
  \text{IM exists} \;\Rightarrow\; \text{Squeezing} \;\Rightarrow\; \text{TLL}
  \;\Rightarrow\; \text{TLL + LDI} \;\Rightarrow\; \text{BV selection (G')}
  \;\Rightarrow\; \text{Regularity (Thm~3.1 of companion paper)}
\]
This is the first system where the entire TLL/LDI $\to$ BV $\to$ regularity
chain closes without open assumptions.
For classical NS ($\beta = 1$), the squeezing property holds but
inertial manifold existence is open---this is precisely the gap
addressed by the $\mu$-reach approach (companion paper).
\end{remark}

% ------------------------------------------------------------------
\subsection{Mathematical tools for TLL and LDI}
% ------------------------------------------------------------------

\begin{remark}[Bridge lemma: squeezing to TLL]
\label{rem:bridge-squeezing}
The squeezing property (Foias--Temam) states that high modes contract faster than low modes diverge. Combined with the determining modes theorem, this implies: $d(u(t), \mathcal{A}) \leq C e^{-\gamma_N t} d(u(0), \mathcal{A}) + C' \sup_{s \leq t} |P_N u(s) - P_N v(s)|$ for some $v \in \mathcal{A}$. Taking logarithms, $\log(1/d(t)) \geq \gamma_N t - C''$, which is the log-modulus bound underlying TLL.
\end{remark}

\begin{remark}[Doubling and Assouad structure]
\label{rem:doubling-assouad}
TLL can be reformulated in terms of the Assouad dimension of the attractor trajectory: if the trajectory $\{u(t) : t \in [0,T]\}$ has Assouad dimension $< 1$, then the ``approximate tangent cone'' has opening angle bounded by $C/|\log r|$, which is equivalent to TLL with rate $1/|\log r|$. This connects the regularity theory to geometric measure theory. Recent structural results on sets of positive reach~\cite{Lytchak2024} provide additional tools for understanding the geometry of attractor neighbourhoods.
\end{remark}

\begin{remark}[LDI as tail-integral equivalence]
\label{rem:ldi-tail}
The LDI condition $\int_0^T \log(1/d(t))\,\mathrm{d}t < \infty$ is equivalent to the tail-integral condition $\int_0^R \frac{1}{s} \left( \int_{\{0 < d \leq s\}} \|u'(t)\|\,\mathrm{d}t \right) \mathrm{d}s < \infty$, relating the integrability of $\log(1/d)$ to the ``speed'' of the trajectory near the attractor. This formulation is more amenable to energy estimates.
\end{remark}

\begin{remark}[Proposition~\ref{prop:variant-B} clarification]
\label{rem:prop-variant-B-clarification}
The BV $\to$ $L^\infty$ embedding in Proposition~\ref{prop:variant-B} requires the additional condition $d(t) > 0$ a.e.: the trajectory must not touch the attractor on a set of positive measure. This is automatically satisfied for Leray--Hopf solutions (which are analytic in space and hence cannot coincide with an attractor element on an interval unless they are that element).
\end{remark}

% ------------------------------------------------------------------
\subsection{Stochastic extensions and advanced tools}
% ------------------------------------------------------------------

\begin{remark}[Stochastic stabilisation]
\label{rem:stochastic-stabilisation}
Adding Feller-type noise $\sigma\,\mathrm{d}W$ to the Navier--Stokes equations can enforce TLL ``almost surely'': the noise prevents the trajectory from approaching the attractor tangentially (which would violate log-Lipschitz), and the Feller property ensures the transition semigroup maps bounded functions to continuous ones, providing the regularity needed for the selection principle.
\end{remark}

\begin{remark}[High-mode stability conjecture]
\label{rem:high-mode-stability}
A sufficient condition for TLL in the NS context: if the high Fourier modes $Q_N u = (I - P_N) u$ satisfy $\|Q_N u(t)\| \leq C N^{-\alpha}$ uniformly on the attractor for some $\alpha > 1/2$, then $d(t) \leq C' N^{-\alpha}$ and TLL follows from the determining modes estimate. This ``high-mode stability conjecture'' is weaker than inertial manifold existence (which requires $\alpha > 1$) and may be verifiable via Galerkin energy estimates. Recent sharp dimension estimates for NS-type attractors~\cite{IlyinKalantarovZelik2025} provide quantitative input for such bounds.
\end{remark}

\begin{remark}[Variational analysis perspective]
\label{rem:variational-analysis}
The theory of metric regularity and variational analysis (Rockafellar--Wets~\cite{RockafellarWets1998}, Dontchev--Rockafellar~\cite{DontchevRockafellar2009}) provides a natural home for the TLL and LDI conditions: TLL is related to the Aubin property of the attractor projection, and LDI is a calmness condition on the distance function. The full toolkit of variational analysis---including Ekeland's variational principle, the Mordukhovich criterion, and the Attouch--Th\'era duality---is applicable and underexplored in this context.
\end{remark}

% ------------------------------------------------------------------
\subsection{Further NS-specific research directions}
% ------------------------------------------------------------------

\begin{remark}[Further research directions]
\label{rem:further-ns-logdist}
\begin{enumerate}
\item \emph{Selection regularity via trajectory attractors:} Formulate the minimal metric derivative selection in the Chepyzhov--Vishik trajectory attractor framework, where non-uniqueness is handled by considering entire trajectories rather than initial data.
\item \emph{Quantitative tracking lemmas:} Establish measure estimates $|\{t : d(t) \leq \varepsilon\}| \leq C \varepsilon^\alpha$ via energy/absorbing-ball arguments and ``no sticky contact'' conditions.
\item \emph{Ekeland variational principle:} Use Ekeland's principle to construct ``almost minimisers'' with controlled variation, bypassing the need for positive reach of $\mathcal{A}$.
\item \emph{Set-valued BV selection:} Determine which dynamical regularity conditions on the flow force bounded Hausdorff variation of the minimiser set $\{u \in \mathcal{A} : d(u, u(t)) = d_{\mathcal{A}}(u(t))\}$.
\item \emph{Approximate prox-regularity:} Establish a modified reach after removing a meager set (Lebesgue-null): the attractor may fail to have positive reach on a negligible set but retain it generically. Results of Lytchak~\cite{Lytchak2024} on the structure of sets of positive reach suggest that such a decomposition may be feasible.
\item \emph{Weak selection functions:} Replace strong Lipschitz projections with selection functions of bounded Hausdorff variation, which exist under weaker regularity assumptions.
\item \emph{Geometric Euler system:} Construct a coherent hierarchy of finite Galerkin approximations $\mathcal{A}_N \to \mathcal{A}$ as a ``reach surrogate'': even if $\mathrm{reach}(\mathcal{A}) = 0$, the Galerkin approximations $\mathcal{A}_N$ have positive reach and converge in Hausdorff distance.
\end{enumerate}
\end{remark}

% ------------------------------------------------------------------
\subsection{Cross-paper connections}
% ------------------------------------------------------------------

\begin{remark}[DFC from turbulence as topological funnel for TLL]
\label{rem:dfc-funnel}
The Downhill Flux Condition (DFC) from the turbulence companion paper~\cite{FST-NS2026} implies that energy flows ``downhill'' through the inertial range. Translated to the attractor geometry, this means trajectories approach $\mathcal{A}$ through a topological funnel: the DFC constrains the approach direction, enforcing a log-Lipschitz rate. This provides a physical mechanism for TLL that is independent of the abstract functional-analytic arguments.
\end{remark}


% ===================================================================
% AI Disclosure
% ===================================================================
\subsection*{AI Disclosure}
In the preparation of this manuscript, AI-assisted tools (Claude, Anthropic)
were used in a supporting capacity, particularly for literature research,
text structuring, and formulation assistance.  The conceptual design,
scientific argumentation, and responsibility for the entire content
rest solely with the author.


\begin{thebibliography}{99}

\bibitem{AmbrosioGigliSavare2008}
L.~Ambrosio, N.~Gigli, and G.~Savar\'e,
\emph{Gradient Flows in Metric Spaces and in the Space of
Probability Measures},
2nd ed., Lectures in Mathematics ETH Z\"urich,
Birkh\"auser, Basel, 2008.

\bibitem{AmbrosioFuscoPallara2000}
L.~Ambrosio, N.~Fusco, and D.~Pallara,
\emph{Functions of Bounded Variation and Free Discontinuity Problems},
Oxford Mathematical Monographs, Clarendon Press, Oxford, 2000.

\bibitem{BabinVishik1992}
A.~V.~Babin and M.~I.~Vishik,
\emph{Attractors of Evolution Equations},
Studies in Mathematics and its Applications, vol.~25,
North-Holland, Amsterdam, 1992.

\bibitem{BaeCannone2016}
H.~Bae and M.~Cannone,
\emph{Log-Lipschitz regularity of the 3D Navier--Stokes equations},
Nonlinear Anal.\ \textbf{137} (2016), 222--245.

\bibitem{Chistyakov2004}
V.~V.~Chistyakov,
\emph{Selections of bounded variation},
J.~Appl.\ Anal.\ \textbf{10} (2004), no.~1, 1--82.

\bibitem{DeBlasiMyjak1993}
F.~S.~De~Blasi and J.~Myjak,
\emph{Ambiguous loci of the nearest point mapping in Banach spaces},
Arch.\ Math.\ \textbf{61} (1993), no.~4, 377--384.

\bibitem{DontchevRockafellar2009}
A.~L.~Dontchev and R.~T.~Rockafellar,
\emph{Implicit Functions and Solution Mappings: A View from
Variational Analysis},
Springer Monographs in Mathematics, Springer, New York, 2009.

\bibitem{ConstantinFoiasTemam1985}
P.~Constantin, C.~Foias, and R.~Temam,
\emph{Attractors representing turbulent flows},
Mem.\ Amer.\ Math.\ Soc.\ \textbf{53} (1985), no.~314.

\bibitem{EdenFoiasNicolaenkoTemam1994}
A.~Eden, C.~Foias, B.~Nicolaenko, and R.~Temam,
\emph{Exponential Attractors for Dissipative Evolution Equations},
Research in Applied Mathematics, vol.~37,
Wiley--Masson, Paris, 1994.

\bibitem{Federer1959}
H.~Federer,
\emph{Curvature measures},
Trans.\ Amer.\ Math.\ Soc.\ \textbf{93} (1959), 418--491.

\bibitem{FoiasManleyRosaTemam2001}
C.~Foias, O.~Manley, R.~Rosa, and R.~Temam,
\emph{Navier--Stokes Equations and Turbulence},
Encyclopedia of Mathematics and its Applications, vol.~83,
Cambridge University Press, 2001.

\bibitem{FST-NS2026}
L.~Geiger,
\emph{A Thermodynamic Obstruction to Finite-Time Blow-Up in the 3D
Navier--Stokes Equations (Conditional Framework)},
Zenodo, 2026.  DOI:\@ to be assigned.

\bibitem{IlyinKalantarovZelik2025}
A.~Ilyin, V.~Kalantarov, and S.~Zelik,
\emph{Attractors and their dimensions for the 3D Fractional
Navier--Stokes--Voigt Equations},
arXiv preprint arXiv:2511.04783, 2025.

\bibitem{Lytchak2024}
A.~Lytchak,
\emph{A note on subsets of positive reach},
Math.\ Nachr.\ \textbf{297} (2024), no.~3, 932--942.

\bibitem{MalletParetSell1988}
J.~Mallet-Paret and G.~R. Sell,
\emph{Inertial manifolds for reaction diffusion equations in higher
space dimensions},
J.~Amer.\ Math.\ Soc.\ \textbf{1} (1988), no.~4, 805--866.

\bibitem{Mane1981}
R.~Man\'e,
\emph{On the dimension of the compact invariant sets of certain
non-linear maps},
Lecture Notes in Mathematics, vol.~898, Springer, Berlin, 1981,
pp.~230--242.

\bibitem{RockafellarWets1998}
R.~T.~Rockafellar and R.~J-B Wets,
\emph{Variational Analysis},
Grundlehren der mathematischen Wissenschaften, vol.~317,
Springer, Berlin, 1998.

\bibitem{PoliquinRockafellarThibault2000}
R.~A.~Poliquin, R.~T.~Rockafellar, and L.~Thibault,
\emph{Local differentiability of distance functions},
Trans.\ Amer.\ Math.\ Soc.\ \textbf{352} (2000), no.~11, 5231--5249.

\bibitem{Saks1937}
S.~Saks,
\emph{Theory of the Integral},
2nd revised ed., Monografie Matematyczne, vol.~7,
G.~E.~Stechert \& Co., New York, 1937.

\bibitem{Temam1997}
R.~Temam,
\emph{Infinite-Dimensional Dynamical Systems in Mechanics and Physics},
2nd ed., Applied Mathematical Sciences, vol.~68,
Springer, New York, 1997.

\bibitem{Zamfirescu1990}
T.~Zamfirescu,
\emph{The nearest point mapping is single valued nearly everywhere},
Arch.\ Math.\ \textbf{54} (1990), no.~6, 563--566.

\bibitem{FoiasTemam1979}
C.~Foias and R.~Temam,
\emph{Some analytic and geometric properties of the solutions of the
evolution Navier--Stokes equations},
J.\ Math.\ Pures Appl.\ \textbf{58} (1979), no.~3, 339--368.

\bibitem{Romanov2000}
A.~V.~Romanov,
\emph{Sharp estimates of the dimension of inertial manifolds for
nonlinear parabolic equations},
Izv.\ Ross.\ Akad.\ Nauk Ser.\ Mat.\ \textbf{57} (1993), 36--52;
English transl.\ in Russian Acad.\ Sci.\ Izv.\ Math.

\end{thebibliography}

\end{document}
